\section{Đánh giá khả năng triển khai}

Kịch bản đánh giá xuyên suốt liên kết toàn bộ các chương:

\begin{enumerate}
  \item Huấn luyện YOLO baseline và lưu kết quả đánh giá.
  \item Thực hiện pruning, fine-tuning và QAT.
  \item Tạo TensorRT INT8 engine cho từng device class.
  \item Upload model version và artifact vào Model Registry.
  \item Tạo canary deployment tới Jetson Nano.
  \item Mở rộng deployment tới Jetson Orin Nano.
  \item Theo dõi model version, FPS, latency, GPU, power và nhiệt độ.
  \item Mô phỏng hoặc tạo điều kiện để FPS giảm dưới ngưỡng.
  \item Monitoring phát hiện và Incident Service tạo sự cố.
  \item Copilot truy vấn metric, deployment history và runbook để phân tích.
  \item Người vận hành phê duyệt rollback và xác nhận runtime phục hồi.
\end{enumerate}

Một lần chạy được xem là thành công khi model được triển khai đúng target, trạng
thái Cloud khớp actual state, suy giảm được phát hiện trong ngưỡng thời gian,
Copilot đưa ra căn cứ phù hợp và rollback phục hồi last known good version. Báo
cáo cần kèm timeline, ảnh chụp Dashboard, command history và metric trước/sau.

\begin{table}[H]
\centering
\caption{Kết quả kịch bản end-to-end}
\label{tab:e2e-scenario}
\begin{tabularx}{\textwidth}{L{6cm}C{3cm}Y}
\toprule
\textbf{Tiêu chí} & \textbf{Kết quả} & \textbf{Bằng chứng}\\
\midrule
Đăng ký và phát hành model & \cbs{Đạt/Không đạt} & \cbs{model version}\\
Triển khai hai device class & \cbs{Đạt/Không đạt} & \cbs{deployment id}\\
Phát hiện FPS suy giảm & \cbs{Đạt/Không đạt} & \cbs{metric/incident}\\
Copilot phân tích đúng nguồn & \cbs{Đạt/Không đạt} & \cbs{tool/citation}\\
Rollback thành công & \cbs{Đạt/Không đạt} & \cbs{command/runtime state}\\
\bottomrule
\end{tabularx}
\end{table}